\section{目录优化：User 与 Core 分离}\label{sec:restructure}

\subsection{问题：用户代码与 ST 代码混在一起}

最初的工程把 \texttt{main.c}（你的应用代码）和 \texttt{system\_stm32f4xx.c}
（ST 的系统代码）都放在 \texttt{Core/src/} 里，\texttt{stm32f4xx\_conf.h}
也放在 \texttt{Core/inc/} 里：

\begin{lstlisting}[style=generic]
Core/
├── src/
│   ├── main.c                 # 用户代码（不该在这里）
│   └── system_stm32f4xx.c     # ST 系统代码
└── inc/
    └── stm32f4xx_conf.h       # 应用配置（不该在这里）
\end{lstlisting}

这样做的坏处是：\textbf{你的代码和厂商代码边界不清}。当升级 CMSIS / 标准外设库
（通常整个替换 \texttt{Driver/} 目录）时，容易误伤自己的代码；反之，
定位「我写的代码」时也要在 ST 源码里翻找。

\subsection{分层后的结构}

将应用层与设备支持层分开：

\begin{lstlisting}[style=generic]
User/                            # 应用层：你写的代码
├── src/main.c                   # 应用入口
└── inc/stm32f4xx_conf.h         # 外设库配置

Core/                            # 设备支持层：ST 系统代码
├── src/system_stm32f4xx.c       # 时钟初始化
└── inc/                         # 预留：设备相关头文件
\end{lstlisting}

分层原则：

\begin{itemize}
  \item \textbf{User/}：业务逻辑、外设应用封装，随项目变化频繁修改；
  \item \textbf{Core/}：与芯片强相关的系统代码，改动少，接近厂商源码；
  \item \textbf{Driver/}：纯第三方库，只读，升级时整体替换。
\end{itemize}

\subsection{对应的 CMake 改动}

源文件列表与头文件路径随之更新：

\begin{lstlisting}[style=generic]
add_executable(${TARGET}
    ${CMAKE_SOURCE_DIR}/User/src/main.c          # 原来在 Core/src
    ${CMAKE_SOURCE_DIR}/Core/src/system_stm32f4xx.c
    ...
)

target_include_directories(${TARGET} PRIVATE
    ${CMAKE_SOURCE_DIR}/User/inc                 # 原来在 Core/inc
    ${CMAKE_SOURCE_DIR}/Core/inc
)
\end{lstlisting}

\subsection{后续可进一步优化的方向}

当前结构已经够用，但如果工程继续增长，可以考虑：

\begin{enumerate}
  \item \textbf{把 Driver 拆成独立仓库 / Git submodule}：
        厂商库体积大且与业务无关，用 submodule 引入可保持主仓库干净、升级方便。
  \item \textbf{引入 \texttt{App/} 或按功能分子目录}：
        当 \texttt{User/src/} 下文件变多时，按功能拆成
        \texttt{User/src/bsp/}、\texttt{User/src/app/} 等，用
        \texttt{target\_sources} 分目录添加。
  \item \textbf{用 CMake 的 object library / interface library}：
        把 \texttt{Driver/}、\texttt{Core/} 各自封装成一个静态库目标，
        主程序只链接这些库，依赖关系更清晰、编译更可复用。
  \item \textbf{用 \texttt{CMakePresets.json} 管理多套配置}：
        例如一份 debug（\texttt{-O0 -g}）一份 release（\texttt{-Os}），
        避免命令行里反复写长参数。
  \item \textbf{提取 CPU 编译选项到工具链文件或一个变量}：
        目前 \texttt{CPU\_FLAGS} 定义在 \texttt{CMakeLists.txt} 顶部，
        若以后支持多个芯片型号，可改成按型号选择的缓存变量。
  \item \textbf{补充 \texttt{.gitignore}}：忽略 \texttt{build/}、
        \texttt{*.elf}/\texttt{*.bin}/\texttt{*.hex}/\texttt{*.map} 等产物，
        目前这些文件仍出现在 \texttt{git status} 里。
\end{enumerate}
